\section{Public Final-Text Predicates and Attack Model}
\label{sec:verification}

The public final-text experiments ask whether a third party, given only text,
can recover a deterministic codeword or at least build a reliable public
predicate for provenance.  The answer in the current artifacts is negative:
the predicates recover zero codeword blocks.  The same predicates nevertheless
have nonzero accept regions, which makes them targets for source-mismatch
optimization.

\paragraph{Public predicate.}
A public final-text predicate is any fixed rule in the current artifact set
that maps final text to a score or accept/reject decision without using
provider-side traces.  The P2 predicates used in the results are shallow
surrogates for final-text observability.  They are not claimed to exhaust all
possible public predicates, and they are not claimed to be deployment-grade
verifiers.  Their diagnostic role is narrower: if they recover no codeword
blocks, they do not verify; if they still accept source-mismatch rows, their
accept region is spoofable.

\paragraph{Attack ladder.}
We evaluate four increasingly strong public-predicate attacks:
\begin{description}
\item[Rejection sampling.] Query the public predicate on existing
source-mismatch rows and select accepts.
\item[Rewrite-lite.] Apply deterministic opener or phrasing edits and requery
the predicate.
\item[Distillation-lite.] Train or fit a lightweight surrogate ranking signal
from public predicate feedback, then rank source-mismatch rows.
\item[Guided rewrite/graft.] Use public predicate feedback to choose and graft
predicate-favorable text fragments onto source-mismatch rows.
\end{description}

These attacks do not prove that every public final-text predicate is broken.
They show that the strongest current final-text surrogate predicates occupy a
bad regime: they do not recover codewords, but their nonzero accept regions can
be found or optimized by source-mismatch examples.

\paragraph{Selection-bias boundary.}
Attack success must not be confused with protected recovery.  The attack
pipeline selects or modifies source-mismatch text using public predicate
feedback.  It does not use the protected trace route, does not recover the
protected codeword, and does not establish natural-output provenance.  Its
purpose is to stress the public final-text predicate substrate.

\paragraph{Naturalness boundary.}
The current guided rewrite/graft attacks are artifact-level predicate attacks.
Their accepted source-mismatch examples should not be described as human-level
natural or semantically equivalent unless a separate naturalness and semantic
preservation evaluation is run.  This limitation weakens any claim about
attack realism, but it does not rescue a predicate that fails codeword recovery
and accepts source-mismatch rows.
